\documentclass[10pt,letterpaper]{article}
\usepackage[utf8]{inputenc}
\usepackage[T1]{fontenc}
\usepackage[margin=0.75in]{geometry}
\usepackage{hyperref}
\usepackage{enumitem}
\usepackage{titlesec}
\usepackage{lmodern}

\pagestyle{empty}

\titleformat{\section}{\large\bfseries\uppercase}{}{0em}{}[\titlerule]
\titlespacing{\section}{0pt}{1ex}{1ex}

\newcommand{\CVItem}[1]{\item #1}

\begin{document}

\begin{center}
    {\LARGE \textbf{Jair Molina Arce}} \\ \vspace{0.5ex}
    Querétaro, Qro., México \\
    \href{mailto:ingjairmolina@gmail.com}{ingjairmolina@gmail.com} | 5652646108 \\
    \href{https://jairmolina.dev}{jairmolina.dev} | \href{https://github.com/smookymolina}{github.com/smookymolina} | \href{https://linkedin.com/in/jair-molina-arce-4909622b2}{linkedin.com/in/jair-molina-arce-4909622b2}
\end{center}

\vspace{1ex}

\begin{center}
    {\Large \textbf{Ingeniero Mecánico}}
\end{center}

\section*{Perfil Profesional}
Ingeniero Mecánico egresado del IPN, actualmente cursando una maestría en Tecnologías Avanzadas. Poseo experiencia comprobada en diseño mecánico, análisis térmico y estructural mediante Elementos Finitos (FEA), dimensionamiento de bancos de prueba e instrumentación de precisión. He caracterizado, diseñado y simulado sistemas mecánicos complejos utilizando SolidWorks, ANSYS y MATLAB/Simulink, y he validado su comportamiento dinámico mediante sistemas de adquisición de datos y pruebas experimentales rigurosas. Mi principal diferencial competitivo es la capacidad de integrar fluidamente el dominio mecánico con la electrónica y la automatización: logro que sensores, actuadores y microcontroladores trabajen en sinergia como un único sistema electromecánico cohesivo. Busco un rol desafiante en ingeniería mecánica orientado al diseño de producto, gestión térmica avanzada, pruebas de confiabilidad o validación de sistemas.

\section*{Habilidades Técnicas}
\begin{itemize}[leftmargin=*, itemsep=0.2ex]
    \CVItem{\textbf{Diseño / CAD-CAE}: SolidWorks (Modelado 3D, Ensamblajes, Dibujos técnicos), ANSYS (Análisis FEA estructural, modal y térmico), MATLAB/Simulink (Modelado dinámico), diseño, cálculo y dimensionamiento de bancos de prueba, lectura e interpretación avanzada de planos bajo normas internacionales (GD\&T), metrología dimensional, impresión 3D (FDM, SLA) para prototipado rápido.}
    \CVItem{\textbf{Térmica y Fluidos}: Análisis de transferencia de calor (conducción, convección, radiación), termodinámica aplicada, gestión térmica integral de gabinetes y equipos electrónicos, cálculo y control de sistemas de ventilación, análisis de eficiencia y balance energético de sistemas.}
    \CVItem{\textbf{Instrumentación y Pruebas}: Adquisición de Datos (DAQ) mediante LabVIEW y Python, selección e implementación de sensores de presión, temperatura (termoacoples, RTDs, NTCs) y flujo, acondicionamiento de señales, estrategias de control PID, calibración de instrumentos, uso de osciloscopio y multímetro, diseño de protocolos de validación experimental y aseguramiento de la trazabilidad metrológica de los ensayos.}
    \CVItem{\textbf{Mecatrónica y Automatización}: Experiencia práctica con microcontroladores (ESP32, STM32, serie RP), programación en C/C++ y Python, control de motores a pasos (NEMA, drivers A4988), puentes H, selección y dimensionamiento de MOSFETs de potencia, buses de comunicación (UART, SPI, I2C), diseño de PCBs de interconexión básica en KiCad.}
    \CVItem{\textbf{Datos y Software}: Programación en Python (pandas, numpy, matplotlib para análisis de datos experimentales), Excel nivel avanzado (Power Query, macros VBA), control de versiones con Git, contenerización básica con Docker.}
\end{itemize}

\section*{Experiencia Relevante}
\begin{itemize}[leftmargin=*, itemsep=0.5ex]
    \item \textbf{Docente Universitario - Ingeniería Térmica y Modelización de Sistemas Aeroespaciales} | \textit{Universidad Internacional de Innovación de Aguascalientes} | Ago. 2025 -- Presente.
    \begin{itemize}[itemsep=0.1ex]
        \item Titular de la materia de Ingeniería Térmica, cubriendo transferencia de calor, termodinámica aplicada y análisis energético de sistemas complejos.
        \item Titular de Modelización de Sistemas Aeroespaciales, abarcando dinámica de vuelo, simulación numérica y desarrollo de modelos matemáticos.
        \item Diseño integral de material didáctico, evaluaciones prácticas basadas en proyectos y simulaciones computacionales, gestionado 100\% en línea.
    \end{itemize}
    \item \textbf{Docente - Tecnologías Disruptivas (Realidad Virtual y Aumentada)} | \textit{ENBA - IPN} | 2024.
    \begin{itemize}[itemsep=0.1ex]
        \item Impartí curso de posgrado enfocado en el impacto y aplicación industrial de la VR/AR.
        \item Diseñé metodologías para la aplicación de estas tecnologías en la capacitación técnica inmersiva y gestión de información espacial.
    \end{itemize}
    \item \textbf{Ingeniería Mecánica e Instrumentación - Banco de Pruebas de Propulsión} | \textit{Instituto Politécnico Nacional (IPN)} | 2022 -- 2024.
    \begin{itemize}[itemsep=0.1ex]
        \item Caractericé, diseñé y simulé la estructura completa de un banco de pruebas mecánicas para cohetes de combustible sólido de investigación.
        \item Ejecuté análisis térmico-estructurales (FEA) de componentes críticos, iterando diseños para soportar cargas dinámicas extremas y gradientes térmicos.
        \item Implementé el sistema DAQ (LabVIEW, Python) logrando una captura de datos de alta fidelidad, vital para la validación de las simulaciones teóricas.
        \item Trabajo de investigación publicado formalmente en \textit{Revista Hacia el Espacio} (2024).
    \end{itemize}
\end{itemize}

\section*{Proyecto Mecánico Destacado}
\begin{itemize}[leftmargin=*, itemsep=0.5ex]
    \item \textbf{Jaguar MX - Sistema de Extracción de Aire Caliente para Gabinete de Telecomunicaciones} (2026).
    \begin{itemize}[itemsep=0.1ex]
        \item \textit{Rol}: Diseño conceptual, gestión térmica, selección de componentes mecánicos/eléctricos y validación de prototipo.
        \item Desarrollé la estrategia completa de disipación térmica del gabinete, estableciendo curvas de operación para extracción de aire caliente.
        \item Calculé flujos de aire requeridos y seleccioné el ventilador industrial MR1238E48B-FSR (48 V) y actuadores de compuerta óptimos.
        \item Instrumenté la medición de temperatura en puntos críticos para garantizar que los equipos de telecomunicaciones operen dentro de su envolvente térmica segura.
        \item Diseñé el layout de los componentes para maximizar el flujo de aire y minimizar la recirculación de aire caliente, modelando escenarios en software.
    \end{itemize}
\end{itemize}

\section*{Proyectos Seleccionados}
\begin{itemize}[leftmargin=*, itemsep=0.2ex]
    \CVItem{\textbf{Impresora 3D DIY - Diseño y Construcción desde Cero}. Diseño íntegro de la estructura mecánica, sistema cinemático CoreXY en SolidWorks, análisis estructural del chasis en aluminio extruido, e integración mecatrónica (Arduino Mega, motores NEMA 17). Cierre del ciclo diseño-manufactura-validación con la impresión exitosa de piezas de prueba dimensionales.}
    \CVItem{\textbf{Pipeline de Automatización de Análisis de Datos y Reportes} (2024 -- 2025). Reducción del tiempo de análisis post-prueba en un 90\% mediante rutinas automáticas que consolidan datos de múltiples sensores en reportes técnicos listos para auditoría.}
    \CVItem{\textbf{Sistema de Control Adaptativo PID con Ajuste Automático} (2023 -- 2024). Modelado matemático de una planta térmica en MATLAB y validación práctica de algoritmos de control.}
\end{itemize}

\section*{Freelance e Investigación}
\begin{itemize}[leftmargin=*, itemsep=0.2ex]
    \CVItem{\textbf{Producción Científica}: Ponente y coautor en el \textit{75th International Astronautical Congress (IAC 2024)}, Milán, Italia (DOI: 10.52202/078365-0120). Publicación en \textit{Revista Hacia el Espacio} (2024). CVU CONACYT: 1340773 | ORCID: 0009-0009-6732-8100.}
    \CVItem{\textbf{Instructor - Excel Avanzado} (Mar. 2026). Automatización de flujos de información técnica y financiera.}
\end{itemize}

\section*{Educación}
\begin{itemize}[leftmargin=*, itemsep=0.2ex]
    \CVItem{\textbf{Maestría en Tecnologías Avanzadas} (2024 -- Presente), Instituto Politécnico Nacional (IPN).}
    \CVItem{\textbf{Ingeniería Mecánica} (2017 -- 2023), Instituto Politécnico Nacional (IPN).}
    \CVItem{\textbf{Bachillerato Técnico en Informática} (2014 -- 2017), Colegio de Bachilleres Plantel 1 "El Rosario".}
\end{itemize}

\section*{Idiomas y Competencias Transversales}
\begin{itemize}[leftmargin=*, itemsep=0.2ex]
    \CVItem{\textbf{Idiomas}: Español (Nativo) | Inglés (Avanzado, con lectura técnica fluida y capacidad de redacción de reportes de ingeniería).}
    \CVItem{\textbf{Soft Skills}: Liderazgo técnico, resolución de problemas complejos, capacidad de análisis crítico, trabajo en equipo multidisciplinario, comunicación efectiva de conceptos técnicos a audiencias no técnicas, adaptabilidad, aprendizaje continuo, gestión del tiempo y organización de proyectos.}
\end{itemize}

\end{document}
